%%%%%%%%%%%%%%%%%%%%%%%%%%%%%%%%%%%%%%%%%%%%%%%%%%%%%%%%%%%%%%%%%%%%%%
% LaTeX Template: Curriculum Vitae
%
% Source: http://www.howtotex.com/
% Feel free to distribute this template, but please keep the
% referal to HowToTeX.com.
% Date: July 2011
%Version for spanish users, by dgarhdez
% 
%%%%%%%%%%%%%%%%%%%%%%%%%%%%%%%%%%%%%%%%%%%%%%%%%%%%%%%%%%%%%%%%%%%%%%
% How to use writeLaTeX: 
%
% You edit the source code here on the left, and the preview on the
% right shows you the result within a few seconds.
%
% Bookmark this page and share the URL with your co-authors. They can
% edit at the same time!
%
% You can upload figures, bibliographies, custom classes and
% styles using the files menu.
%
% If you're new to LaTeX, the wikibook is a great place to start:
% http://en.wikibooks.org/wiki/LaTeX
%
%%%%%%%%%%%%%%%%%%%%%%%%%%%%%%%%%%%%%%%%%%%%%%%%%%%%%%%%%%%%%%%%%%%%%%
\documentclass[paper=a4,fontsize=11pt]{scrartcl} % KOMA-article class
							
\usepackage{hyperref}
\usepackage[english]{babel}
\usepackage[utf8x]{inputenc}
\usepackage[protrusion=true,expansion=true]{microtype}
\usepackage{amsmath,amsfonts,amsthm}     % Math packages
\usepackage{graphicx}                    % Enable pdflatex
\usepackage[svgnames]{xcolor}            % Colors by their 'svgnames'
\usepackage{geometry}
	\textheight=700px                    % Saving trees ;-)
\usepackage{url}

\frenchspacing              % Better looking spacings after periods
\pagestyle{empty}           % No pagenumbers/headers/footers

%%% Custom sectioning (sectsty package)
%%% ------------------------------------------------------------
\usepackage{sectsty}

\sectionfont{%			            % Change font of \section command
	\usefont{OT1}{phv}{b}{n}%		% bch-b-n: CharterBT-Bold font
	\sectionrule{0pt}{0pt}{-5pt}{1pt}}

%%% Macros
%%% ------------------------------------------------------------
\newlength{\spacebox}
\settowidth{\spacebox}{8888888888}			% Box to align text
\newcommand{\sepspace}{\vspace*{1em}}		% Vertical space macro

\newcommand{\MyName}[1]{ % Name
                \Huge \hfill #1
                \par \normalsize \normalfont}
		
\newcommand{\MySlogan}[1]{ % Slogan (optional)
		\large \usefont\hfill \textit{#1}
		\par \normalsize \normalfont}

\newcommand{\NewPart}[1]{\section*{{#1}}}

\newcommand{\PersonalEntry}[2]{
		\noindent\hangindent=2em\hangafter=0 % Indentation
		\parbox{\spacebox}{        % Box to align text
		\textit{#1}}		       % Entry name (birth, address, etc.)
		\hspace{1.5em} #2 \par}    % Entry value

\newcommand{\SkillsEntry}[2]{      % Same as \PersonalEntry
		\noindent\hangindent=2em\hangafter=0 % Indentation
		\parbox{\spacebox}{        % Box to align text
		\textit{#1}}			   % Entry name (birth, address, etc.)
		\hspace{1.5em} #2 \par}    % Entry value	
		
\newcommand{\EducationEntry}[4]{
		\noindent \textbf{#1} \hfill      % Study
		\colorbox{White}{%
			\parbox{5cm}{%
			\hfill\color{Black}#2}} \par  % Duration
		\noindent \textit{#3} \par        % School
		\noindent\hangindent=2em\hangafter=0 \small #4 % Description
		\normalsize \par}

\newcommand{\WorkEntry}[4]{				  % Same as \EducationEntry
		\noindent \textbf{#1} \hfill      % Jobname
		\noindent\colorbox{Black}{\color{White}#2} \par  % Duration
		\noindent \textit{#3} \par              % Company
		\noindent\hangindent=2em\hangafter=0 \small #4 % Description
		\normalsize \par}

%%% Begin Document
%%% ------------------------------------------------------------
\begin{document}
% you can upload a photo and include it here...
%\begin{wrapfigure}{l}{0.5\textwidth}
%	\vspace*{-2em}
%		\includegraphics[width=0.15\textwidth]{photo}
%\end{wrapfigure}

\MyName{Michelle Gomes Pereira}



%%% Personal details
%%% ------------------------------------------------------------
\NewPart{Dados pessoais}{}
\begin{minipage}{0.55\textwidth}
	\begin{tabular}{|p{\textwidth}}
		\textbf{Nascimento:} 18/09/1987 \newline
		\textbf{Estado civil:} Casado \newline
		\textbf{Email:} \href{mailto:mmichelli6@gmail.com}{mmichelli6@gmail.com} \newline
		\textbf{Redes Sociais:} \newline
		\href{https://www.linkedin.com/in/michellegomespereira/}{LinkedIn }
		\href{https://www.facebook.com/michelle.gomes.904}{ Facebook}
	\end{tabular}
\end{minipage}%
\begin{minipage}{0.45\textwidth}
	\begin{tabular}{|p{\textwidth}}
		\textbf{Rua:} Acácio Costa Junior, 153 \newline
		\textbf{Bairro:} Juliana \newline
		\textbf{Cep:} 31744-490 - Belo Horizonte/MG \newline
		\textbf{Tel.:} +55 31 985650762 \newline
	\end{tabular}
\end{minipage}%


%%% Dados Pessoais
%%% ------------------------------------------------------------
\NewPart{Resumo}{}
Sou Biologa com interesse na área de \textit{Bioinformática}, \textit{Data Science} e \textit{Big Data}. Cursei disciplinas da 
pósgraduação no DCC - Departamento de Ciências da Computação, tais como: Ambientes de Computação(Python), Modularidade POO (Java), 
Bioinformática e Empreendedorismo. Utilizei ferramentas tais como: VSCode, Git, GitHub, RStudio, R, Python, Pipenv, Jupyter, Java, Eclipse, 
MS Teams, Miro, Google Meet, Google Colab.


%%% Formação Acadêmica
%%% ------------------------------------------------------------
\NewPart{Formação Acadêmica}{}

\EducationEntry{Puntifícia Universidade Católica}{2022 - 2023}{Pós-graduação (especialização)}{
	\begin{itemize}
		\item{Data Science e Big Data }

	\end{itemize}
}

\EducationEntry{Universidade Federal de Minas Gerais}{2013 - 2020}{Graduação}{
	\begin{itemize}
		\item{Ciências Biológicas }

	\end{itemize}
}

%%% Experiência Profissional
%%% ------------------------------------------------------------
\NewPart{Experiência Profissional}{}

\EducationEntry{Analista Desenvolvedor Sr – Invillia.}{}{Outubro/2020 até Dezembro/2020}{
	\begin{itemize}
		\item{Desenvolvimento de projetos em squad de criação e monitoramento de microserviços
		      utilizando a linguagem Kotlin/Java no backend. Criação de proxy para publicação de
		      microserviços com o Apigee. Seguindo princípios da arquitetura hexagonal.
			  Criação de testes unitários e de integração alocado no C6bank. Utilizou-se as 
			  seguintes especificações e tecnologias: \newline \newline 
		      \textit{Kotlin, Java, Backend For Frontend, Spring, Cloud AWS, Koin, Javalin, Exposed, Fuel, Swagger, OpenApi, WireMock, MockK, REST-
			      assured, H2, flywaydb, Gradle, Vault, Concorce CI, Rundeck CD, Ansible, Bitbucket,
			      Sonar, Artifactory, Graylog, Grafana, Prometheus, App Dynamics, Apigee, Docker, Splunk,
			      Zabbix, Fluentd, Harbor Registry, Trivy Vulnerability Scanner.}}

	\end{itemize}
}

\EducationEntry{Analista Desenvolvedor Sr – Pif Paf Alimentos S.A.}{}{Março/2019 até Outrubro/2020}{
	\begin{itemize}
		\item{Desenvolvimento e implementação de sistemas utilizando a plataforma Java como backend
		e o javascript(node) como front-end e práticas ágeis juntamente com a cultura de DevOps.
		Automatização de pipeline utilizando Gitlab e automatização de testes e criação de métricas
		e configuração do Sonarqube. Integrações de sistemas utilizando princípios de clean architecture 
		e a especificação microprofile para integrações com sistemas SAP e gateway de pagamento. 
		Utilizou-se as seguintes especificações e tecnologias: \newline \newline
		      \textit{SAP Hybris Commerce, Spring Boot, Spring Framework, Spring Data, Spring Security, Spring MVC,
			  Microprofile, Quarkus, Cloud AWS, Docker, Minikube, Helm, Ansible, Terraform, Python, Gradle,
			  Jquery, node, gulp, mocha, chai, javascript, Bootstrap, Intellij, Eclipse,VSCode, Postman, SoapUI, 
			  Sonarqube, GitLab, Git, Sourcetree, MySQL, VPN, ImpEx, Flexible Search.}}

	\end{itemize}
}

\EducationEntry{Analista Desenvolvedor Sr – FÓTON Informática S.A}{}{Setembro/2018 até Março/2019}{
	\begin{itemize}
		\item{Desenvolvimento e implementação de sistemas utilizando práticas ágeis com scrum e as
		especificações Java Enterprise Edition. Integrações de sistemas utilizando o estilo
		arquitetural REST e SOAP para integração com sistemas, realização de testes
		automatizados. Utilizou-se as seguintes especificações e tecnologias: \newline \newline
		      \textit{Java EE, JBoss EAP 6.1, JSF, JAX-WS, JAX-RS, SQL, JPA, JMS, Primefaces, Jquery,
			  Arquillian, Eclipse IDE, Docker, Maven, Sonatype Nexus, Jasper Studio Professional, Postman,
			  SoapUI, GitLab, Git, Sourcetree, TFS, PostgreSQL, PGAdmin}}

	\end{itemize}
}

\EducationEntry{Analista Desenvolvedor – SQUADRA TECNOLOGIA}{}{Novembro/2017 até Setembro/2018}{
	\begin{itemize}
		\item{Desenvolvimento e implementação de sistemas utilizando as especificações Java
		Enterprise Edition. Integrações de sistemas utilizando o estilo arquitetural REST para
		integração com sistemas Bancários e WebService SOAP para integrações com ERP Totvs
		(DataSul). Utilizou-se as seguintes especificações e tecnologias: \newline \newline
		      \textit{Java EE, JBoss EAP 7.1, JSF 2.2, JAX-WS 2.0, JAX-RS 2.0, OAuth 2, Quartz Job
			  Scheduling, Struts 2, Docker, Intersystem Caché, SQL, ORM Baseado em Refleção proprietário -
			  SmartMVCP, Primefaces, Jquery, Velocity, Eclipse IDE, Android Studio, Android SDK,
			  Jasper Studio Professional, Postman, SoapUI, Bugzilla.}}

	\end{itemize}
}

\EducationEntry{Analista Desenvolvedor – AXXIOM}{}{Agosto/2017 até Novembro/2017}{
	\begin{itemize}
		\item{Desenvolvimento de projeto de SmartGrid da rede de distribuição da compainha energética
		de Minas Gerais utilizando as especificações Java Enterprise Edition e a plataforma
		Android alocado na CEMIG. Utilizou-se as seguintes especificações e tecnologias: \newline \newline
		      \textit{Java EE, Docker, JSF, EJB, JPA, CDI, Primefaces, Hibernate, Wildfly, Git, IBM Informix, SQL,
			  Eclipse IDE, Maven, Android Studio, Gradle, Android SDK, SoapUI, Postman.}}

	\end{itemize}
}

\EducationEntry{Analista Desenvolvedor – SQUADRA TECNOLOGIA}{}{Dezembro/2016 até Agosto/2017}{
	\begin{itemize}
		\item{Desenvolvimento e implementação de sistemas de logística de coleta de BAGs e Amostras
		laboratoriais alocado no Instituto Hermes Pardini (Processo de negócio consistia em partes
		Web e partes Mobile), sistema de integração (WebService) com sistema bancário via API
		RESTful e integração com ERP Totvs V12 via o protocolo SOAP para emissão e quitação
		(baixa) de títulos via EDI - Electronic Data Interchange. Utilizou-se as seguintes
		especificações e tecnologias: \newline \newline
		      \textit{Java EE 7, Docker, JSF 2.2, JAX-WS 2.0, JAX-RS 2.0, OAuth 2, Quartz Job Scheduling, Struts 2,
			  Intersystem Caché, SQL, ORM Baseado em Refleção proprietário - SmartMVCP,
			  Primefaces, Jquery, Velocity, Eclipse IDE, Android Studio, Android SDK, Jasper Studio
			  Professional, Postman, SoapUI.}}

	\end{itemize}
}

\EducationEntry{Analista Desenvolvedor Java – Montreal Informática SA}{}{Setembro/2016 até Outubro/2016}{
	\begin{itemize}
		\item{Participação em projeto de migração de servidor de aplicação alocado na JUCEMG.
		Realizando testes e alterações ao consumir WebServices com certificados digitais e
		correções de Bugs. Os projetos utilizavam o JBoss AS 6 e foram migrados para o Wildfly
		10. Utilizou-se as seguintes especificações e tecnologias: \newline \newline
		      \textit{Java EE 7, JSF 2.2, EJB 3.1, JPA 2.1, CDI 1.1, JAX-WS 2.0, Keycloack, RichFaces,
			  JQuery, Hibernate, Wildfly 10, SVN Subversion, Oracle Database, SQL, Eclipse IDE,
			  SoapUI, VPNs.}}

	\end{itemize}
}

\EducationEntry{Analista de Sistemas – Maxinst Consultoria e Tecnologia Ltda}{}{Setembro/2015 até Janeiro/2016}{
	\begin{itemize}
		\item{Atribuições / Serviços Realizados: Atendimento a chamados dos clientes, desenvolvimento
		e aplicação de melhorias ao sistema IBM Máximo. Utilizou-se as seguinte ferramentas e
		tecnologias: \newline \newline
		      \textit{WebLogic, WebSphere, IBM DB2, SourceTree, Git, SQL, Java, Eclipse IDE, IBM
			  Máximo.}}

	\end{itemize}
}

\EducationEntry{Programador Java – Basis Tecnologia da Informação SA}{}{Janeiro/2014 até Setembro/2015}{
	\begin{itemize}
		\item{Atribuições / Serviços Realizados: Desenvolvimento de sistemas utilizando práticas de
		desenvolvimento ágil e testes unitários e de integração automatizados, utilizando
		integração contínua e as seguintes ferramentas e tecnologias: \newline \newline
		      \textit{Plataforma Java, SQL, JIRA, Astah Community, SVN, Git, OracleSQLDeveloper,
			  Stash, Maven, Jenkins, SonarQube, iReport Designer, Eclipse IDE, SourceTree -
			  Atlassian, Wildfly 8.2,GWT, SmartGWT, JSON, HTML, CSS, SmartClient, Nexus, Docker,
			  Banco de dados Oracle 11g, SQL Server, Flywaydb, Arquillian, JUnit, Web Server, Spring Framework,
			  PiketLink, EJB, Hibernate JPA, CDI, JSF, Primefaces, VPNs}}

	\end{itemize}
}

\EducationEntry{Analista Desenvolvedor Java – FIEMG}{}{Setembro/2010 até Agosto/2013}{
	\begin{itemize}
		\item{Desenvolvimento, implementação e implantação de sistemas utilizando práticas ágeis sobre 
		a plataforma Java no STI - Setor de Inovação Tecnológica do CETEC/FIEMG. Utilizando as seguintes ferramentas e tecnologias: \newline \newline
		      \textit{JEE, Eclipse IDE, JBoss Application Server, Apache Tomcat, Google Web Toolkit,
			  MySQL, MySQL Workbench, PostgreSQL, XWiki Enterprise, XWiki Enterprise
			  Manager, Servidor Linux CentOS, Protocolo de comunicação Cliente/Servidor SSH,
			  Linux, Elicitação de Requisitos, UML - Unified Modeling Language, Astah
			  Community, CMapTools - Knowledge Modeling Kit, FreeMind - Mind mapping
			  software, Dia Diagram Editor, MS Project, WBS Chart Pro.}}

	\end{itemize}
}

\EducationEntry{Estagiário Prog. PHP – Fundação Centro Tecnológico de Minas Gerais}{}{Janeiro/2010 até Agosto/2010}{
	\begin{itemize}
		\item{Participação no desenvolvimento do SIGLa, um sistema
		de gerenciamento integrado de laboratórios, utilizando as seguintes tecnologias: \newline \newline
		      \textit{Controle de versionamento SVN, Zend Framework (MVC), Doctrine (ORM), Dojo
			  toolkit(View), Desenvolvimento em ambiente Linux Ubuntu 7.10, Apache HTTP
			  Server, SGBD - MySql 5, Ferramenta de gerenciamento do Banco de Dados –
			  PHPMyAdmin, Linguagem de programação PHP OO (Orientado a Objetos), IDE
			  Eclipse PDT}}

	\end{itemize}
}

\EducationEntry{Técnico em Eletrônica – Fundação Centro Tecnológico de Minas Gerais}{}{Março/2004 até Agosto/2008}{
	\begin{itemize}
		\item{Modernização e operacionalização de um reator PVD (Phisical Vapor Deposition).
		Participação na implementação de circuitos eletrônicos de interface com módulos 
		de aquisição de dados (DAQ / National Instruments) para controle e monitoramento 
		através de Software (LabView / National Instruments) e sistemas supervisorios. \newline \newline
		      \textit{Atuando nas seguintes áreas:
			  Programação PLC, Linguagem C, Microcontroladores, Sistemas Eletrônicos de Medida e de Controle de Processos, Circuitos Eletrônicos Analógicos e Digitais, Instrumentação Eletrônica, Dispositivos de Potência, Transdutores, Atuadores.}}

	\end{itemize}
}


%%% Certificação
%%% ------------------------------------------------------------
\NewPart{Certificações}{}
\EducationEntry{ORACLE CERTIFIED ASSOCIATE, JAVA SE 8 PROGRAMMER}{}{CertiProf - Setembro/2018}{}
\EducationEntry{DEVOPS ESSENTIALS PROFESSIONAL CERTIFICATE - DEPC}{}{CertiProf - Outubro/2018}{}
\EducationEntry{SCRUM FOUNDATION PROFESSIONAL CERTIFICATE - SFPC}{}{CertiProf - Novembro/2018}{}
\EducationEntry{KANBAN FOUNDATION}{}{CertiProf - Dezembro/2018}{}


%%% Work experience
%%% ------------------------------------------------------------
\NewPart{Cursos Extra Curricular}{}

\EducationEntry{Scalable Microservices with Kubernetes}{Junho/2021}{by Google - Udacity}{}
\EducationEntry{Julia Scientific Programming}{Abril/2021}{University of Cape Town}{}
\EducationEntry{R Programming}{Março/2021}{Jonhs Hopkins University}{}
\EducationEntry{Python and Flask Framework}{Fevereiro/2021}{Udemy}{}
\EducationEntry{Deep Learning with Python}{Janeiro/2021}{Udemy}{}
\EducationEntry{Python and Data Science for beginners}{Dezembro/2020}{Udemy}{}
\EducationEntry{Orquestração de Conteiners com Kubernetes}{Janeiro/2018}{Udemy}{}
\EducationEntry{Developing Android Apps}{Janeiro/2017}{Udacity}{}
\EducationEntry{Android Development for Beginners}{Novembro/2016}{Udacity}{}
\EducationEntry{Desenvolvimento de Aplicações JEE}{Outubro/2013}{Framework System}{}
\EducationEntry{Redes de computadores}{Fevereiro/2005}{Training 2020}{}
\EducationEntry{Gerência de projetos}{Julho/2004}{CETEC}{}

%%% Skills
%%% ------------------------------------------------------------
\NewPart{Idiomas}{}

\SkillsEntry{Idiomas}{Português (fluente)}
\SkillsEntry{}{Inglês (intermediário)}
\SkillsEntry{}{Espanhou (básico)\\}
\SkillsEntry{Proficiência}{Inglês - CENEX FALE UFMG }

%\SkillsEntry{}{}

%%% Work experience
%%% ------------------------------------------------------------
\NewPart{Palestra, Seminário e Workshop}{}

\EducationEntry{Deep Dive: Oracle MySQL Cloud Service}{Novembro/2016}{Speaker: Morgan Tocker,
	MySQL Product Manager, ORACLE}{}
\EducationEntry{Primefaces Responsivo}{Novembro/2016}{AlgaWorks}{}
\EducationEntry{Oficina de Srping Framework}{Agosto/2016}{AlgaWorks}{}
\EducationEntry{Dev App com RAD Studio XE2, DataSnap e
	Firemonkey para Windows, Mac OS, iPhone e iPad}{Maio/2012}{Delphi Meeting 2012}{}
\EducationEntry{Hardwares de Aquisição de Dados com LabVIEW}{Abril/2010}{NI National Instruments do Brasil}{}
\EducationEntry{Software livre e novos paradigma}{Março/2005}{SENAI}{}
\EducationEntry{Arquitetura de sistemas Automatizados}{Março/2005}{SENAI}{}
\EducationEntry{LabView (NI) “Painel Virtual do Reator BAS”}{Outubro/2004}{CETEC}{}
\EducationEntry{Automação dos processos industriais}{Março/2004}{SENAI}{}

%%% References
%%% ------------------------------------------------------------
%\NewPart{References}{}
%Available upon request
\end{document}